\metatitle{Stable polynomials}

\metadescription{Multivariate stable polynomials, strongly Rayleigh measures,
stability-preserving operations, and same-phase stability.}
\metakeywords{stable polynomial, real stable polynomial, strongly Rayleigh
measure, Borcea--Branden symbol, multiaffine polynomial, same-phase stability}


\section[stablePolynomial]{Stable polynomials}



Let $\mathcal{H} \subset \setC$ denote the
upper half-plane $\{z \in \setC: \mathrm{im}(z) > 0 \}$.

\begin{definition}[Stability]

A multivariate polynomial $P \in \setC[z_1,\dotsc,z_n]$
is called \defin{stable}\label{stablePolynomialDefinition}
if it does not vanish on $\mathcal{H}^n$.
That is, $P$ is stable if
\[
 \zvec^* \in \mathcal{H}^n \implies P(\zvec^*) \neq 0.
\]
One can easily show that $P \in \setR[z_1,\dotsc,z_n]$ is stable if and only
if $P(\alpha+\lambda t)=0$ has only real zeros for all
$\alpha \in \setR^n$, $\lambda \in \setR_+^n$.
A univariate polynomial with real coefficients
is stable if and only if all its roots are real.
\end{definition}

The Lee--Yang circle theorem is a foundational circular-domain analogue.
For a finite ferromagnetic Ising model in a uniform external field, every zero
of the partition function, viewed as a polynomial in the fugacity, lies on the
unit circle \cite{LeeYang1952}.  A Möbius transformation between the disk and
the upper half-plane transports such zero-free regions into the language of
stability; this viewpoint underlies the modern Lee--Yang program
\cite{BorceaBranden2009I}.





\begin{definition}[Strongly Rayleigh, \cite{BorceaBrandenLiggett2008}]
\label{stronglyRayleigh}
A measure $\mu$ on the set of subsets of $[n]$ is called
\defin{strongly Rayleigh} if
\[
 \sum_{S \subseteq [n]} \mu(S) \prod_{s\in S} x_s
\]
is real stable.
\end{definition}


\subsection[stabilityOperators]{Operators preserving stability}

\begin{theorem}[Stability preserving operators]
The following operations on polynomials preserve stability.
Let $p,q \in \setR[z_1,\dotsc,z_n]$ be stable polynomials.
Then the following are also stable.
\begin{itemize}
\item $p\cdot q$ (Product)
\item $\partial_{z_j} p$ (Differentiation)
\item $z_j \partial_{z_j} p$ (Degree-preserving differentiation)
\item $p(z_1,\dotsc,z_{j-1},r,z_{j+1},\dotsc,z_n)$ (Real specialization)
\item $p(z_1,\dotsc,z_{j-1},z_1,z_{j+1},\dotsc,z_n)$ (Projection)
\item $z_j^d p(z_1,\dotsc,z_{j-1},-1/z_{j},z_{j+1},\dotsc,z_n)$ (Inversion)
\end{itemize}
Here, $d$ is the degree of $z_j$ in $p$.
\end{theorem}


\subsection[borceaBrandenSymbols]{Borcea--Brändén symbols}

The Borcea--Brändén classification packages a linear operator into a
polynomial, or into an entire power series, called its
\defin{symbol}\label{stabilitySymbol}.  This is often the most practical way
to prove that an operator preserves stability: instead of testing the operator
on all stable inputs, one tests one universal object.  For real stability one
must allow both signs.  For instance,
\(f(z)\mapsto f(-z)\) preserves real-rootedness, but it is detected by the
minus symbol rather than the plus symbol.

\subsubsection[algebraicStabilitySymbols]{Algebraic symbols}

\begin{definition}[Algebraic symbols]
Let $\mathbf{d}=(d_1,\dotsc,d_n) \in \setN^n$, and let
$\setR_{\leq \mathbf{d}}[z_1,\dotsc,z_n]$ denote the polynomials whose degree
in $z_j$ is at most $d_j$.
For a linear operator
\[
 T:\setR_{\leq \mathbf{d}}[z_1,\dotsc,z_n]\to \setR[x_1,\dotsc,x_m],
\]
define the two \defin{algebraic symbols}\label{algebraicSymbol} of $T$ by
\[
 G_T^{+}(\xvec,\mathbf{w})
 \coloneqq
 T\left( \prod_{j=1}^n (z_j+w_j)^{d_j} \right),
 \qquad
 G_T^{-}(\xvec,\mathbf{w})
 \coloneqq
 T\left( \prod_{j=1}^n (z_j-w_j)^{d_j} \right).
\]
Here $T$ acts on the $z$-variables and the $w_j$'s are auxiliary variables.
\end{definition}

\begin{theorem}[Borcea--Brändén finite-degree classification, see \cite{BorceaBranden2009I}]
\label{borceaBrandenFiniteSymbol}
Let \(T:\setR_{\leq \mathbf{d}}[z_1,\dotsc,z_n]\to
\setR[x_1,\dotsc,x_m]\) be linear.  Then \(T\) sends every real stable
polynomial to a real stable polynomial or to zero if and only if one of the
following holds.
\begin{itemize}
\item \(G_T^{+}\) is real stable.
\item \(G_T^{-}\) is real stable.
\item The range of \(T\) has dimension at most two and can be written as
  \(T(f)=\alpha(f)P+\beta(f)Q\), where \(\alpha,\beta\) are real linear
  functionals and \(P,Q\) are real stable polynomials in proper position.
  In one common convention this means that \(Q+iP\) is stable.
\end{itemize}
\end{theorem}

For operators over complex coefficients, the corresponding stability-preserver
classification has only the plus symbol, apart from the rank-one exceptional
case.  The extra sign in the real theorem records the distinction between
preserving the upper half-plane and preserving real-rootedness.

\begin{example*}[Two elementary symbols]
For \(T=\partial_{z_i}\), the plus symbol is
\[
  G_T^{+}(\xvec,\mathbf{w})
  =
  d_i(x_i+w_i)^{d_i-1}
  \prod_{j\neq i}(x_j+w_j)^{d_j},
\]
which is stable.  This recovers the stability-preserving property of
partial differentiation.

For the univariate reflection \(T(f)(x)=f(-x)\) on
\(\setR_{\leq d}[z]\), the plus symbol \((w-x)^d\) is not stable, but the
minus symbol
\[
  G_T^{-}(x,w)=(-x-w)^d
\]
is stable.  Hence the sign alternative is genuinely needed for real
stability preservers.
\end{example*}

\subsubsection[transcendentalStabilitySymbols]{Transcendental symbols}

\begin{definition}[Laguerre--Pólya class and transcendental symbol]
Let \(\mathcal{LP}_r\) be the
\defin{Laguerre--Pólya class}\label{laguerrePolyaClass}
of entire functions in \(r\) variables
that are limits, uniformly on compact sets, of real stable polynomials.
This is the multivariate Laguerre--Pólya class.

For a linear operator
\[
  T:\setR[z_1,\dotsc,z_n]\to \setR[x_1,\dotsc,x_m],
\]
its \defin{transcendental symbol}\label{transcendentalSymbol}
is the formal power series
\[
  \overline{G}_T(\xvec,\mathbf{w})
  \coloneqq
  T\!\left(e^{-\zvec\cdot\mathbf{w}}\right)
  =
  \sum_{\alpha\in\setN^n}
    (-1)^{|\alpha|}
    T(\zvec^\alpha)
    \frac{\mathbf{w}^{\alpha}}{\alpha!}.
\]
\end{definition}

\begin{theorem}[Borcea--Brändén transcendental classification, see \cite{BorceaBranden2009I}]
\label{borceaBrandenTranscendentalSymbol}
Up to the same rank-at-most-two exception as above, a linear operator
\[
  T:\setR[z_1,\dotsc,z_n]\to \setR[x_1,\dotsc,x_m]
\]
preserves real stability if and only if
\[
  \overline{G}_T(\xvec,\mathbf{w})\in \mathcal{LP}_{m+n}
  \qquad\text{or}\qquad
  \overline{G}_T(\xvec,-\mathbf{w})\in \mathcal{LP}_{m+n}.
\]
\end{theorem}

This is the infinite-degree analogue of the algebraic-symbol test.  It also
explains why the classical \hyperref[multiplierSequences]{multiplier
sequence} theorem has an exponential generating function.  If
\[
  T_\gamma(z^k)=\gamma_k z^k,
\]
then
\[
  \overline{G}_{T_\gamma}(x,w)
  =
  \sum_{k\geq 0}\gamma_k\frac{(-xw)^k}{k!}
  =
  \Phi_\gamma(-xw).
\]
Thus the Pólya--Schur condition on \(\Phi_\gamma\) is the diagonal-operator
case of the Borcea--Brändén symbol criterion.  Compare this with
\hyperref[polyaFrequency]{Pólya frequency sequences}: PF sequences are
ordinary coefficient sequences controlled by Toeplitz
\hyperref[totalPositivity]{total positivity},
whereas multiplier sequences are diagonal operators controlled by exponential
symbols.

\name{Ming-Jian Ding} and \name{Bao-Xuan Zhu} give a systematic treatment of
stability for combinatorial polynomials \cite{DingZhu2021Stability}.  They
apply Borcea--Br\"and\'en type characterizations to real and Hurwitz stability,
Tur\'an expressions, alternating-runs polynomials, and several interlacing and
alternatingly-increasing phenomena.

\subsection[gardingPolynomials]{G{\aa}rding polynomials and Rayleigh applications}

\name{Hao Fang} and \name{Biao Ma} introduce
\defin{G{\aa}rding polynomials}\label{gardingPolynomial}, a class strictly
containing real stable
polynomials \cite{FangMa2026x}.  The definition is phrased in terms of
positivity regions invariant under translation by positive vectors; they prove
equivalent formulations by polarization and by a recursive criterion using
partial derivatives.  Multi-affine G{\aa}rding polynomials with nonnegative
coefficients satisfy the Rayleigh property, and their positive univariate
specializations have ultra-log-concave coefficient sequences.

Their matroid applications cover cycle matroids of series--parallel
networks, \hyperref[uniformMatroid]{uniform matroids}, matroids obtained from
a uniform matroid by removing one basis, and every matroid on at most six
elements.  For these classes, the spanning-set, cospanning-set, and basis
generating polynomials are G{\aa}rding.  All of these polynomials, as well as
the corresponding independent-set generating polynomials, are Rayleigh, and
their univariate specializations satisfy ultra-log-concavity and Mason-type
inequalities \cite[Thm.~1.3]{FangMa2026x}.


\subsection[gurvitsCapacity]{Capacity and coefficient bounds}

For a homogeneous polynomial $p\in\setR[x_1,\dotsc,x_n]$ of degree $n$ with
nonnegative coefficients, define its \defin{capacity}\label{stableCapacity} by
\[
  \operatorname{Cap}(p)
  \coloneqq
  \inf_{x_1,\dotsc,x_n>0}\frac{p(x_1,\dotsc,x_n)}{x_1\dotsm x_n}.
\]

\begin{theorem}[Gurvits, \cite[Thm.~2.4 and Cor.~2.5]{Gurvits2008}]
If $p$ is real stable, then
\[
  \frac{\partial^n p}{\partial x_1\dotsm\partial x_n}(0)
  \geq \frac{n!}{n^n}\operatorname{Cap}(p).
\]
\end{theorem}

This turns stability into a coefficient lower bound and behaves well under
successive differentiation.  For a doubly stochastic matrix
$A=(a_{ij})$, the polynomial
\[
  p_A(\xvec)=\prod_{i=1}^n\left(\sum_{j=1}^n a_{ij}x_j\right)
\]
is real stable and has capacity one, while the coefficient of
$x_1\dotsm x_n$ is $\operatorname{per}(A)$.  The theorem therefore gives the
van der Waerden bound $\operatorname{per}(A)\geq n!/n^n$.



\subsection[stabilitymultiaffine]{Multi-affine stability}

A polynomial $P(z_1,\dotsc,z_n)$ is
\defin{multi-affine}\label{multiAffinePolynomial}
if the degree in each variable is at most $1$.



\begin{theorem}[See \cite{Branden2007}]
A homogeneous multi-affine (real) polynomial $f(x_1,\dotsc,x_n)$ is real
stable if and only if
\[
 \frac{\partial f}{\partial x_i} \cdot \frac{\partial f}{\partial x_j} -
 \frac{\partial^2 f}{\partial x_i \partial x_j} \cdot f \geq 0
\]
whenever $x_1,x_2,\dotsc,x_n \in \setR$.
\end{theorem}
A polynomial satisfying this inequality is said to have the
\defin{strongly Rayleigh property}\label{stronglyRayleighProperty}.
The \defin{weak Rayleigh property}\label{weakRayleighProperty}
holds if the above inequality is true whenever the variables are positive
real numbers.


\begin{theorem}[Grace--Walsh--Szeg\H{o}, see \cite{Grace1902,Szego1922,Walsh1922,RahmanSchmeisser2002}]
Let $P(z_1,\dotsc,z_n)$ be a \emph{symmetric} multi-affine polynomial.
Then $P$ is stable if and only if the univariate diagonal specialization
$P(t,\dotsc,t)$ is stable.
If $P$ has real coefficients, this is equivalent to the univariate
specialization $P(t,\dotsc,t)$ being real-rooted.
\end{theorem}

\begin{definition}[Apolar polynomials]
Let
\[
  F(z)=\sum_{k=0}^n \binom{n}{k} a_k z^k
  \qquad\text{and}\qquad
  G(z)=\sum_{k=0}^n \binom{n}{k} b_k z^k.
\]
The polynomials $F$ and $G$ are \defin{apolar}\label{apolarPolynomial} if
\[
  \sum_{k=0}^n (-1)^k \binom{n}{k} a_k b_{n-k}=0.
\]
\end{definition}

\begin{theorem}[Grace's apolarity theorem, see \cite{Grace1902,Szego1922,RahmanSchmeisser2002}]
Let $C$ be a circular domain, and let $F$ and $G$ be apolar
polynomials of degree $n$.  If all zeros of $F$ lie in $C$, then $G$
has at least one zero in $C$.
\end{theorem}

A useful way to reduce questions to the multi-affine case is
\defin{polarization}\label{polarizationDefinition}.
If $0 \leq k \leq d_j$, set
\[
 \operatorname{Pol}_{d_j}(z_j^k)
 \coloneqq \binom{d_j}{k}^{-1}
 \elementaryE_k(z_{j,1},\dotsc,z_{j,d_j}).
\]
For monomials, define
\[
 \operatorname{Pol}_{\mathbf{d}}(\zvec^\alpha)
 \coloneqq
 \prod_{j=1}^n \operatorname{Pol}_{d_j}(z_j^{\alpha_j}),
\]
and extend linearly to polynomials.  This gives the polarization
$\operatorname{Pol}_{\mathbf{d}}(P)$ of a polynomial whose degree in $z_j$
is at most $d_j$.

\begin{theorem}[Polarization, see \cite{BorceaBranden2009II}]
Let $P \in \setC[z_1,\dotsc,z_n]$ have degree at most $d_j$ in the variable
$z_j$.  Then $P$ is stable if and only if
$\operatorname{Pol}_{\mathbf{d}}(P)$ is stable.
\end{theorem}


\section[stabilityConnections]{Jump systems and matroids}


For $P = \sum_{\alpha} c_\alpha \zvec^\alpha$, let
$\supp(P) \coloneqq \{\alpha : c_\alpha \neq 0\}$.

\begin{theorem}[See \cite{Branden2007}]
If $P$ is stable, then $\supp(P)$ is a jump system.
\end{theorem}

For multiaffine polynomials, stability controls more than the underlying
support.  Suppose
\[
  f(\zvec)=\sum_{S\subseteq[n]}a_S\zvec^S
\]
is real stable and $a_\emptyset>0$.  Then the coefficient signs
$S\mapsto\sign(a_S)$ orient the $\Delta$-matroid $\supp(f)$
\cite[Thm.~3.1]{Chin2026x}.  The construction is compatible with changing the
reference feasible set through stability-preserving inversions.

There is also a tropical analogue.  If the coefficients $c_S$ lie in the
complex Puiseux series and $\sum_Sc_S\zvec^S$ is multiaffine and stable, then
\[
  S\longmapsto\operatorname{val}(c_S)
\]
is a valuated $\Delta$-matroid \cite[Thm.~4.1]{Chin2026x}.  Here valuated
$\Delta$-matroid means that every cell in the induced regular subdivision is
a $\Delta$-matroid polytope.  The supports need not satisfy strong exchange,
so this conclusion does not automatically give the stronger constant-parity
version of valuated $\Delta$-matroids.  See also the
\hyperref[deltaMatroid]{matroid page}.


\begin{theorem}[See \cite{ChoeOxleySokalWagner2004}]
If $P$ is homogeneous and multi-affine stable,
then its support is the set of bases of a \hyperref[matroid]{matroid}.
See also \cite[Cor. 3.4]{Branden2007}.
\end{theorem}
The converse is in general not true but some classes of \hyperref[matroid]{matroids}
always have a \hyperref[basisGeneratingPolynomial]{basis-generating polynomial}
which is stable.
For example, \hyperref[matroidRegular]{regular matroids}
(which include \hyperref[matroidGraphic]{graphic matroids}),
\hyperref[matroidUniform]{uniform matroids} and
a subclass of \hyperref[transversalMatroids]{transversal matroids}
(see \cite{ChoeOxleySokalWagner2004}).


\hyperref[matroidOperations]{Deletion and contraction} on matroids preserve
stability of basis-generating polynomials.

For a connected graph $G$ on vertex set $V$, the spanning-tree degree
enumerator
\[
  \sum_T \prod_{v \in V} x_v^{\deg_T(v)},
\]
where the sum is over all spanning trees $T$ of $G$, is real stable if and
only if $G$ is distance-hereditary \cite{CherkashinPetrovProzorov2023}.


\begin{theorem}[See \cite{BrandenHuh2020}]
Let $P \in \setR[z_1,\dotsc,z_n]$ be homogeneous, real stable, and have
nonnegative coefficients.  Then $P$ is \hyperref[lorentzianPolynomial]{Lorentzian}.
\end{theorem}


\subsection[determinantalStability]{Determinantal stability}

\begin{theorem}[Borcea--Brändén,
\cite[Prop.~1.12]{BorceaBranden2010}; see also
\cite[Prop.~2.4]{BorceaBranden2008MixedDeterminants}]
Let $A,B_1,\dotsc,B_n$ be Hermitian $m\times m$ matrices, and suppose that
$B_1,\dotsc,B_n$ are positive semidefinite.  Then
\[
 \det(A+z_1B_1+\dotsb+z_nB_n)
\]
is stable, in fact real stable, or identically zero.
\end{theorem}

For real values of the variables, the matrix pencil is Hermitian and its
determinant is real, so the determinant polynomial has real coefficients.
The proof rests on the interaction between the upper half-plane and positive
semidefiniteness.  Suppose that all $\operatorname{Im}(z_i)>0$ and that
$M=A+\sum_i z_iB_i$ is singular.  For a nonzero vector $v$ in its kernel,
\[
 0=\operatorname{Im}(v^*Mv)
  =\sum_i \operatorname{Im}(z_i)v^*B_iv.
\]
Every summand is nonnegative, so $B_iv=0$ for every $i$, and then $Av=0$.
Thus $v$ lies in the kernel of every matrix in the pencil, making the
determinant identically zero.  Hence a nonzero determinant of this form cannot
vanish when all variables lie in the upper half-plane.


\subsection[mixedCharacteristicPolynomial]{Mixed characteristic polynomials and barriers}

For positive semidefinite Hermitian matrices $A_1,\dotsc,A_m$, define
\[
 \mu[A_1,\dotsc,A_m](x)
 \coloneqq
 \left.
 \prod_{i=1}^m(1-\partial_{z_i})
 \det\left(xI+\sum_{i=1}^m z_iA_i\right)
 \right|_{z_1=\dotsb=z_m=0}.
\]
This is the \defin{mixed characteristic polynomial}
\label{mixedCharacteristicPolynomialDefinition}.

\begin{theorem}[Marcus--Spielman--Srivastava,
\cite[Cor.~4.4 and Thm.~5.1]{MarcusSpielmanSrivastava2015II}]
\label{mixedCharacteristicPolynomialBound}
The polynomial $\mu[A_1,\dotsc,A_m](x)$ is real-rooted.  If additionally
\[
  \sum_{i=1}^m A_i=I
  \qquad\text{and}\qquad
  \operatorname{tr}(A_i)\leq\epsilon\quad\text{for every $i$},
\]
then its largest zero is at most $(1+\sqrt{\epsilon})^2$.
\end{theorem}

Real-rootedness follows by applying the stability-preserving operators
$1-\partial_{z_i}$ to the determinantal polynomial.  The root bound is proved
by tracking logarithmic-derivative barrier functions as these operators are
applied.  Together with \hyperref[interlacingFamilies]{interlacing
families}, this is the main polynomial tool in the solution of the
Kadison--Singer problem.


\subsection[bivariateStableDeterminantalRepresentation]{Bivariate determinantal representations}

\begin{theorem}[Borcea--Brändén,
\cite[Thm.~1.13 and Cor.~6.7]{BorceaBranden2010}]
Let $P(z,w)\in\setR[z,w]$ have degree $d$.  Then $P$ is real stable if and
only if there are real symmetric $d\times d$ matrices $A,B,C$, with $B$ and
$C$ positive semidefinite, such that
\[
 P(z,w)=\pm\det(A+zB+wC).
\]
\end{theorem}

This converse is the two-variable stable-polynomial form of the classical Lax
theorem.  In its ternary form, the theorem states that if a homogeneous
$h(x,y,t)\in\setR[x,y,t]$ of degree $d$ is hyperbolic with respect to
$(1,0,0)$, then
\[
 h(x,y,t)=h(1,0,0)\det(xI+yB+tC)
\]
for real symmetric $d\times d$ matrices $B,C$.  The Lax conjecture was proved
by \name{J. William Helton} and \name{Victor Vinnikov}, and by
\name{Adrian Lewis}, \name{Pablo Parrilo}, and \name{Motakuri Ramana}
\cite{HeltonVinnikov2007,LewisParriloRamana2005}.

To obtain the bivariate converse, one homogenizes $P$ to
\[
 P_H(z,w,t)\coloneqq t^dP(z/t,w/t).
\]
Stability makes $P_H$ hyperbolic with respect to every direction
$(a,b,0)$ with $a,b>0$.  Applying the ternary Lax theorem and using all these
directions forces the matrices multiplying $z$ and $w$ to be positive
semidefinite, which gives the representation above.

The bivariate conclusion does not extend to a polynomial-level statement in
arbitrarily many variables.  \name{Petter Brändén} used the stable
basis-generating polynomial $h_{V_8}$ of the Vámos matroid to construct the
real-zero polynomial $p(\xvec)=h_{V_8}(\mathbf{1}+\xvec)$, for which no
positive power $p^N$ has a determinantal representation
\cite[Thm.~3.3]{Branden2011}.  Consequently, $h_{V_8}$ and its powers need not
themselves have the analogous homogeneous determinant representation.  This
does not disprove the generalized cone-level Lax conjecture: that conjecture
asks only whether every hyperbolicity cone is spectrahedral, possibly through
a different defining polynomial, and remains open.


\subsection[combinatorialStablePolynomialExamples]{Combinatorial stability examples}

\begin{example*}[The monotone column permanent conjecture]

Let $A$ be a real $n\times n$-matrix in which entries along columns decrease.
Let $Z_n = \mathrm{diag}(z_1,\dotsc,z_n)$ be a diagonal matrix with indeterminates,
and let $J_n$ be the matrix with all entries equal to $1$.
Then
\[
 \mathrm{per}(J_nZ_n + A) = \mathrm{per}\left( (z_j + a_{ij})_{1 \leq i,j \leq n} \right)
\]
is stable, see \cite{BrandenHaglundVisontaiWagner2011}.
By noting that the \hyperref[eulerianPolynomial]{Eulerian polynomials}
are generated by excedances,
they give a new proof that (a multivariate version of) Eulerian polynomials
are stable.
\end{example*}


\begin{example*}[Multivariate Eulerian polynomials]
\label{ex:stableMultivariateEulerianOperator}
\name{Petter Brändén} (see \cite[Thm. 2.5]{HaglundVisontai2012}) proved
stability of the following multivariate refinement of Eulerian polynomials.
Set
\[
 \widetilde A_n(\xvec,\yvec)
 =
 \sum_{\pi \in \symS_n}
 \prod_{\substack{1\leq i<n\\ \pi_i \gt \pi_{i+1}}} x_{\pi_i}
 \prod_{\substack{1\leq i<n\\ \pi_i \lt \pi_{i+1}}} y_{\pi_{i+1}}.
\]
Thus descent tops receive \(x\)-variables and ascent tops receive
\(y\)-variables.  The proof uses a stability-preserving operator.  If
\[
  \partial
  =
  \sum_{i=1}^n \frac{\partial}{\partial x_i}
  +
  \sum_{i=1}^n \frac{\partial}{\partial y_i},
\]
then
\[
  T_{n+1}
  \coloneqq
  (x_{n+1}+y_{n+1}) + x_{n+1}y_{n+1}\partial
\]
preserves real stability on multi-affine polynomials, by the
\hyperref[borceaBrandenFiniteSymbol]{Borcea--Brändén symbol criterion}.
Equivalently, if
\[
  U=\prod_{i=1}^n (x_i+u_i)(y_i+v_i),
\]
then the algebraic symbol
\[
  T_{n+1}(U)
  =
  \left(
    x_{n+1}+y_{n+1}
    +x_{n+1}y_{n+1}
    \sum_{i=1}^n
    \left(\frac{1}{x_i+u_i}+\frac{1}{y_i+v_i}\right)
  \right)U
\]
is stable.
The recurrence
\[
  \widetilde A_{n+1}(\xvec,\yvec)
  =
  T_{n+1}\widetilde A_n(\xvec,\yvec)
\]
then proves stability by induction from \(\widetilde A_1=1\).
Specializing \(y_i=1\) and then \(x_i=t\) gives
\[
  t\,\widetilde A_n(t,\dotsc,t;1,\dotsc,1)=A_n(t),
\]
so the \hyperref[eulerianZeros]{classical Eulerian polynomial} has only real
roots.

In fact, \name{Jim Haglund} and \name{Mirko Visontai} also treat Stirling
permutations, where plateaus occur \cite{HaglundVisontai2012}.  If \(Q_n\) is
the set of Stirling permutations and \(D(\sigma),A(\sigma),P(\sigma)\) are the
descent-top, ascent-top, and plateau sets, then
\[
 C_n(\xvec,\yvec,\zvec)
 =
 \sum_{\sigma \in Q_n}
 \prod_{i\in D(\sigma)} x_{\sigma_i}
 \prod_{i\in A(\sigma)} y_{\sigma_i}
 \prod_{i\in P(\sigma)} z_{\sigma_i}
\]
is stable.  Specializing \(y_i=z_i=1\) recovers real-rootedness of the
second-order Eulerian polynomials.
\end{example*}




\subsection[stableMatchingNegativeDependence]{Matching polynomials and negative dependence}

\begin{example*}[Multivariate matching polynomials]
The multivariate \hyperref[realRootedMatching]{matching polynomial} is real
stable; see \cite{BorceaBranden2009I,BorceaBranden2009II}.
\name{Jonathan D. Leake} and \name{Nick R. Ryder} distinguish two types of
multivariate matching polynomials, recording either vertex labels or edge
labels in the matchings \cite{LeakeRyder2019}:
\begin{align}
\mu_V(G;x_1,\dotsc,x_n)
  & \coloneqq
  \sum_{\substack{M \subseteq E(G) \\ \text{$M$ matching}}}
  \prod_{uv \in M} -x_u x_v, \\
\mu_E(G;x_1,\dotsc,x_n)
  & \coloneqq
  \sum_{\substack{M \subseteq E(G) \\ \text{$M$ matching}}}
  \prod_{e \in M} x_e.
\end{align}
\end{example*}

The \defin{multi-affine vertex matching polynomial}
\label{multiAffineVertexMatchingPolynomial}
$\mu_V(G;x_1,\dotsc,x_n)$ is real stable, see \cite{BorceaBranden2009I}.

\name[D.G. Wagner]{David Wagner}'s survey \cite{Wagner2009StableSurvey}
is a useful reference for the theory and applications of multivariate
stable polynomials, including the Borcea--Brändén classification of
stability preservers, matrix applications, and probabilistic applications.

\begin{theorem}[Borcea--Brändén--Liggett,
\cite[Thm.~4.9]{BorceaBrandenLiggett2008}]
Let $\mu$ be a probability measure on subsets of $[n]$ whose generating
polynomial
\[
  \sum_{S\subseteq[n]}\mu(S)\prod_{i\in S}x_i
\]
is real stable.  Then every measure obtained from $\mu$ by applying
positive external fields is negatively associated: increasing functions
depending on disjoint coordinate sets have nonpositive covariance.
\end{theorem}

In fact the conclusion is the stronger property CNA+: it persists after
conditioning, projection, and positive external fields.  Determinantal
measures induced by positive contractions are strongly Rayleigh; this includes
product measures and uniform random spanning-tree measures
\cite[Prop.~3.5]{BorceaBrandenLiggett2008}.




\subsection[samePhaseStable]{Same phase stability}

\name{Jonathan D. Leake} and \name{Nick R. Ryder} introduced the following
notion, which is strictly weaker than stability \cite{LeakeRyder2019}.

\begin{definition}[Same-phase stability]
 A polynomial $P(z_1,\dotsc,z_n) \in \setR[z_1,\dotsc,z_n]$ is said
 to be \defin{same-phase stable}\label{samePhaseStabilityDefinition}
 if for every $\lambda \in \setR_+^n$, the univariate polynomial
 $P(\lambda_1 t,\lambda_2 t,\dotsc,\lambda_n t) \in \setR[t]$ is real-rooted.
\end{definition}

A family of polynomials $\{p_j(\xvec) \}_{j=1}^m$ is
\defin{same-phase compatible}\label{samePhaseCompatible}
if for every vector of positive numbers $\lambda$, the univariate polynomials
$\{p_j(\lambda t) \}_{j=1}^m$ are
\hyperref[compatiblePolynomials]{compatible}.


Let $P, P_0, \dotsc, P_m$ be polynomials such that
$P = P_0 + \sum_{j=1}^m z_{i_j} P_{j}$.
This is a \defin{proper splitting of $P$}\label{properSplitting}
if none of the polynomials $P_j$ depends on the variables
$\{z_{i_1}, \dotsc, z_{i_m} \}$.


\begin{theorem}[See \cite{LeakeRyder2019}]
Suppose $P$ is a multi-affine polynomial with nonnegative coefficients.
Then the following are equivalent:
\begin{itemize}
\item $P$ is same-phase stable;
\item For any proper splitting,
$P = P_0 + \sum_{j=1}^m z_{i_j} P_{j}$,
the polynomials $P_0, z_{i_1} P_{1}, \dotsc, z_{i_m} P_{m}$
are same-phase compatible.

\item Some proper splitting
$P = P_0 + \sum_{j=1}^m z_{i_j} P_{j}$
such that $P_0, z_{i_1} P_{1}, \dotsc, z_{i_m} P_{m}$ are same-phase
compatible.

\end{itemize}

\end{theorem}


The \defin{multi-affine edge matching polynomial}
\label{multiAffineEdgeMatchingPolynomial}
$\mu_E(G;x_1,\dotsc,x_n)$ is same-phase stable \cite[Cor. 3.2]{LeakeRyder2019}.

Let $G$ be a graph on the vertex set $[n]$, and define the
\defin{multivariate independence polynomial}\label{multivariateIndependencePolynomial}
\[
I(G;x_1,\dotsc,x_n)
\coloneqq
\sum_{\substack{A \subseteq V(G) \\ A \text{ independent}}}
\prod_{v \in A} x_v.
\]
Leake and Ryder prove that $I(G;x_1,\dotsc,x_n)$ is same-phase stable
\emph{if and only if} $G$ is \hyperref[clawFreeGraph]{claw-free}
\cite{LeakeRyder2019}.
This is a reinterpretation of Engström's weighted result \cite{Engstrom2007},
but Leake and Ryder give a much sleeker proof, and the if-and-only-if
condition explains
that the claw-free requirement is sharp.


\begin{example*}[Multivariate Eulerian polynomials II]
In joint work with \name{Olivia Nabawanda}, we show that the following
generalization of the \hyperref[eulerianPolynomial]{Eulerian polynomial} is
same-phase stable, but not stable \cite{AlexanderssonNabawanda2021}:
\[
 \tilde{A}_n(\xvec) = \sum_{\pi \in \symS_n} \prod_{\pi_i \gt \pi_{i+1}} x_{i}.
\]
Note that we now keep track of the descent index, rather than the descent
bottom value.
Moreover, we showed that
\[
 \tilde{A}_{n-1}(\lambda_1 t, \dotsc, \lambda_{n-1} t)
 \interl
 \tilde{A}_n(\lambda_1 t, \dotsc, \lambda_n t)
\]
for all positive real numbers $\lambda_1,\dotsc,\lambda_n$.
\end{example*}
